\chapter{Lezione 17/03}

Il motivo per cui introduciamo questo concetto di sostituzione è perché il lemma o teorema che poi introdurremo legato a questo concetto di sostituzione ci permetterà di ottenere formule valide a partire da altre formule valide.

Questo sarà una prima forma di deduzione logica molto elementare, ma che ci introdurrà verso calcoli deduttivi più raffinati e complessi.

\begin{theorem}[namedlabel={thm:substitution}{Teorema di Sostituzione}]{Teorema di Sostituzione}
  Sia \(\phi\in TAUT\) e \(\sigma\) un'arbitraria sostituzione. Allora \(\phi\sigma\in TAUT\).
\end{theorem}

\begin{note}{}
  Una formulazione alternativa e più intuitiva di questo teorema è che la sostituzione preserva la tautologicità (o la validità) delle formule.
\end{note}

Per dimostrare questo teorema, andremo a introdurre una versione semantica della sostituzione.

Dati \(\phi, \sigma\) e \(\alpha\) (assegnamento per \(\phi\sigma\)), definiamo \(\alpha^\sigma:AP(\phi)\to\set{0,1}\) come segue:
\[\alpha^\sigma(P)=val^\alpha(\sigma(P))\]

Possiamo vedere l'applicazione di \(\sigma\) come una mera sostituzione sintattica (renaming), mentre questa nuova definizione ci dà una sostituziuone semantica.

Vediamo dunque un risultato intermedio che ci sarà utile per dimostrare il \Namedref{thm:substitution}.
\begin{lemma}[namedlabel={lem:substitution}{Lemma di Sostituzione}]{Lemma di Sostituzione}
  \(\forall \phi \in PL, \sigma \in PL^{AP}, \alpha \in ASS(\phi\sigma)\) vale che
  \[val^\alpha(\phi\sigma)=val^{\alpha^\sigma}(\phi)\]
\end{lemma}


\begin{proof}
  Dimostriamo per induzione strutturale.

  Nel caso base \(\phi \in AP\). Allora \(\phi\sigma = \sigma(\phi)\).

  Ma quindi
  \[val^{\alpha}(\phi\sigma) = val^\alpha(\sigma(\phi)) = \alpha^\sigma(\phi) \overset{def}{=} val^{\alpha^\sigma}(\phi)\]

  Andando quindi ai casi induttivi, iniziamo con \(\phi = \neg \psi\). In questo caso dalla definizione di sostituzione abbiamo
  \[\phi\sigma = (\neg\psi)\sigma = \neg(\psi\sigma)\]

  Allora:
  \[val^{\alpha}(\phi\sigma)= val^{\alpha}(\neg(\psi\sigma)) = 1 - val^{\alpha}(\psi\sigma)\]
  Ma per ipotesi induttiva \(val^{\alpha}(\psi\sigma) = val^{\alpha^\sigma}(\psi)\), quindi
  \[val^{\alpha}(\phi\sigma)= 1 - val^{\alpha^\sigma}(\psi) = val^{\alpha^\sigma}(\neg\psi) = val^{\alpha^\sigma}(\phi)\]

  Per gli altri casi induttivi la struttura di dimostrazione è analoga, e si basa sempre sull'ipotesi induttiva.
\end{proof}

Andiamo quindi alla dimostrazione del \Namedref{thm:substitution}.
\begin{proof}{\Namedref{thm:substitution}}
  Vogliamo dimostrare che \(\forall \sigma \in PL^{AP}\), \(\phi \in TAUT \implies \phi\sigma \in TAUT\).

  Procediamo per contrapposizione, ovvero assumiamo che \(\phi\sigma \not\in TAUT\) e dimostriamo che \(\phi \not\in TAUT\).

  Se \(\phi\sigma \not\in TAUT\) allora \(\exists \alpha \in Ass(\phi\sigma)\st val^\alpha(\phi\sigma) = 0\).

  Ma per il \Namedref{lem:substitution}, \(val^\alpha(\phi\sigma) = val^{\alpha^\sigma}(\phi)\), quindi \(val^{\alpha^\sigma}(\phi) = 0\), e quindi \(\phi \not\in TAUT\).
\end{proof}

\begin{observation}{}
  È importante notare che nonostante la sostituzione preserva validità e tautologicità, essa non preserva la soddisfacibilità o la verità in generale.

  Formalmente:
  \[\phi \in Sat \not\implies \phi\sigma \in Sat\]

  Questo è facile da vedere considerando \(\phi = P \in AP\). Si ha che \(\phi \in Sat\) banalmente da qualsiasi assegnamento \(\alpha\) tale che \(\alpha(P)=1\).
  Scegliendo però \(\sigma : P \to P\land \neg P\), che è chiaramente insoddisfacibile.
  Di conseguenza, fissato un modello \(m\), si ha che
  \[m\models\phi \not\implies m\models\phi\sigma\]
  Altrimenti otterei una contraddizione.

  Cosa diversa per la insoddisfacibilità. Infatti
  \[\phi \in Unsat \iff \neg \phi \in TAUT \implies (\neg\phi)\in TAUT\iff \]
  \[\iff \neg(\phi\sigma) \in TAUT \implies \phi\sigma \in Unsat\]
\end{observation}

Possiamo notare che \(A \to (B\to C)\equiv (A\land B)\to C\). Chiamiamo la prima formula \(\phi_1\) e la seconda \(\phi_2\).

Sappiamo che \(\phi_1\) può essere falsa se e soltanto se \(A\) è vera, \(B\) è vera e \(C\) è falsa. 
Invece essendo \(\phi_2\) una semplice implicazione questa è falsa se e soltanto se \(A\land B\) è vera e \(C\) è falsa, ovvero se e soltanto se \(A\) è vera, \(B\) è vera e \(C\) è falsa. Quindi \(\phi_1\) è falsa se e soltanto se \(\phi_2\) è falsa, e quindi \(\phi_1\equiv\phi_2\).

Possiamo trasformare questa equivalenza in una tautologia, scrivendo \(A \to (B\to C) \leftrightarrow (A\land B)\to C\).

\begin{theorem}[namedlabel={thm:semantic_deduction}{Teorema di Deduzione Semantica}]{Teorema di Deduzione Semantica}
  Sia \(\Gamma \subseteq PL\) e \(\phi,\psi \in PL\). Allora
  \[\Gamma,\phi \models \psi \iff \Gamma \models \phi\to\psi\]
\end{theorem}

\begin{proof}
  Dalla definizione di conseguenza logica, \(\Gamma,\phi \models \psi\) significa che
  \[\forall m \in 2^{AP}, m\models \Gamma \cup \set{\phi} \implies m\models \phi\]

  L'argomento sinistro dell'implicazione è però equivalente a dire che \(m\models \Gamma\) e \(m\models \phi\). Quindi
  \[\forall m \in 2^{AP}, m\models \Gamma \text{e} m\models \phi \implies m\models \psi\]

  Dando quindi semantica agli operatori di metalivello nel modo usuale (ovvero identico a quelli della logica proposizionale), e possiamo riscrivere la precedente formula con variabili proposizionali \(A\), \(B\) e \(C\) al posto di \(m\models \Gamma\), \(m\models \phi\) e \(m\models \psi\) rispettivamente, ottenendo
  \[\forall m \in 2^{AP}, A\text{ e }B \implies C\]

  Ma dalla tautologia vista in precedenza, abbiamo che
  \[\forall m \in 2^{AP}, m\models \Gamma \implies (m\models \phi \implies m\models \psi)\]

  Per semantica dell'implicazione però \(m\models \phi \implies m\models \psi \iff m\models \phi \to \psi\), dunque
  \[\forall m \in 2^{AP}, m\models \Gamma \implies m\models \phi\to\psi\]
\end{proof}

\section{Forme normali}

Abbiamo un linguaggio che contiene molti simboli, ma abbiamo già osservato tramite alcune tatutologie che è possibile esprimere alcuni operatori in funzione di altri.

Un esempio potrebbero essere le leggi di De Morgan
\[\neg(\phi\land\psi) \equiv \neg\phi \lor \neg\psi\]
\[\neg(\phi\lor\psi) \equiv \neg\phi \land \neg\psi\]

Interpretando queste equivalenze come regole di trasformazione, questo ci permette di spingere le negazioni fino alle proposizioni atomiche mantenedo l'equivalenza logica.

In questo modo, possiamo definire una prima forma normale, chiamata PNF (Positive Normal Form) o NNF (Negation Normal Form), che contiene tutte le formule di PL in cui le negazioni occorrono solo davanti a proposizioni atomiche.

Per definire formalmente questa forma normale, introduciamo il concetto di letterale

\begin{definition}{Letterale}
  Una formula \(\phi\) è un letterale se \(\phi \in AP\) o \(\phi = \neg P \), con \(P \in AP\).
\end{definition}

\begin{definition}{Forma normale positiva}
  Definiamo la forma normale positiva (PNF o NNF) come l'insieme di tutte le formule costruite come in PL, senza usare negazioni e utilizzando l'insieme dei letterali come AP
\end{definition}

Si ha ovviamente che \(PNF \subset PL\), ma dalle leggi di De Morgan è ovvio che \(\forall \phi \in PL \exists \phi' \in PNF: \phi \iff \phi'\)

In generale il concetto di forma normale è un sott'insieme di formule costruibile in maniera induttiva tale che mantenga completezza semantica, ovvero che per ogni formula esista una formula equivalente in forma normale.

Le due forme normali più importanti sono la forma normale congiuntiva (CNF) e la forma normale disgiuntiva (DNF).

\begin{definition}{Disjuntiva Normal Form}
  Definiamo la forma normale disgiuntiva (DNF) come l'insieme di tutte le formule con la seguente struttura:
  \[\phi = \bigvee_{i=1}^k(\bigwedge_{j=1}^{r_i} l_{ij})\]
  dove \(l_{ij}\) è un letterale per ogni \(i\) e \(j\).
\end{definition}

\begin{example}{}
  Un esempio di formula in DNF è la seguente:
  \[(P\land \neg Q) \lor (\neg P \land \neg Q \land R) \lor \ldots\]

  Ognuna delle parentesi (quindi ogni congiunzione) è detta clausola congiuntiva.
\end{example}

\begin{definition}{Conjunctive Normal Form}
  Definiamo la forma normale congiuntiva (CNF) come l'insieme di tutte le formule con la seguente struttura:
  \[\phi = \bigwedge_{i=1}^k(\bigvee_{j=1}^{r_i} l_{ij})\]
  dove \(l_{ij}\) è un letterale per ogni \(i\) e \(j\).
\end{definition}

\begin{example}{}
  Un esempio di formula in CNF è la seguente:
  \[(P\lor \neg Q) \land (\neg P \lor \neg Q \lor R) \land \ldots\]

  Ognuna delle parentesi (quindi ogni disgiunzione) è detta clausola disgiuntiva.
\end{example}

Essendoci dualità tra congiunzione e disgiunzione, è facile vedere che le due forme normali sono duali.

Ovviamente \(CNF \cup DNF \subseteq PNF\).

\todo{Esempio tabella verità.}

\chapter{Lezione 19/03}

Per la \(\CNF\) il ragionamento è simile, poiché dobbiamo esprimere la formula come una congiunzione di clausole disgiuntive, che cozza con l'uso degli assegnamenti come in precedenza, poiché gli assegnamenti sono delle congiunzioni.

Quindi invece di costruire una formula elencando tutti i casi in cui è vera ma costruiremo una formula che non è vera in tutti gli assegnamenti in cui quella originale era falsa.

Quindi basta creare una formula in cui ogni clausola è della forma forma analoga a quella della \(\DNF\) ma con una negazione davanti, tutte queste congiunte. Ma dalla leggi di de morgan queste clausole congiuntive sono trasformabili in clausole disgiuntive.

\todo{fai esempio pratico}

Se sono tutti a \(1\) posso prendere una qualsiasi tautologia che è già in \(\CNF\) come \(P \lor \neg P\) e mappare in questa.

Queste costruizioni non sono di certo ottimali, poiché per ogni \(\phi\) la formula \(\phi'\) sarà molto più grande. Infatti la grandezza della tabella di verità è \(2^{\AP{\phi}}\), e quindi nel peggiore dei casi si può avere un numero esponenziale sulle proposizioni atomiche di clausole. Non è detto che sia sempre così ma ci sono alcune formule in cui è inevitabile, in particolare quando si passa dall'una all'altra.

\begin{example}{}
  Sia \(\phi \in\DNF\) tale che \(\phi\) sia:
  \[(P_1 \land P_2)\lor (P_3 \land P_4)\]

  Andiamo quindi a trasformarla in \(\CNF\) usando le tautologie note
  \[(P_1 \land P_2)\lor (P_3 \land P_4)\equiv\]
  \[\equiv ((P_1 \land P_2)\lor P_3)\land((P_1 \lor P_4)\lor P_4)\equiv\]
  \[\equiv (P_1 \lor P_3) \land (P_2 \lor P_3) \land (P_1 \lor P_4) \land (P_2 \lor P_4)\]

  Si è passati quindi da \(2\) a \(4\), che può sembrare un incremento lineare. Possiamo notare però che questo processo può essere vista come una funzione di scelta.

  In particolare, data
  \[\phi = (l_1^1 \land l_2^1) \lor (l_1^2 \land l_2^2) \lor\ldots\lor (l_1^n \land l_2^n)\]
  Possiamo costruire una classe di funzione \(f: \set{1,\ldots,n}\to \set{1,2}\) tale che ogni clausola della trasformazione corrisponde a una di queste funzioni di scelta, dunque avremo \(\abs{\phi'}\leq 2^{k}\).

  SPIEGA MEGLIO
\end{example}

Il motivo per cui abbiamo introdotto queste forme è perché se la formula è in \(\DNF\) sarà molto computazionalmente facile valutare al sua soddisfacibilità, mentre se è in \(\CNF\) sarà molto facile valutare la sua validità.

Infatti sia \(\phi = C_1 \lor C_2 \lor\ldots\lor C_k\), con \(C_i = (l_1 \land \ldots \land l_n), \forall i \in \set{1,\ldots, k}\). Una formula del genere è soddisfacibile se e solo se almeno una delle \(C_i\) è soddisfacibile, ma ogni \(C_i\) è insoddisfacibile se e solo se esistono due letterali \(l_j\) e \(l_{j'}\) tali che \(l_j = \neg l_{j'}\). Questo perché se ci sono tali letterali, la clausola contiene la contraddizione \(P \land \neg P\), altrimenti è sembre possibile costruire un assegnamento che rende vera la clausola, in maniera analoga ma inversa a come abbiamo ragionato precedentemene con le tabelle di verità.

Dualmentne, se \(\phi = D_1 \land D_2 \land\ldots\land D_k\), con \(D_i = (l_1 \lor \ldots \lor l_n), \forall i \in \set{1,\ldots, k}\). Una formula del genere è valida se e solo se tutte le \(D_i\) sono valide, ma ogni \(D_i\) è valida se e solo se esistono due letterali \(l_j\) e \(l_{j'}\) tali che \(l_j = \neg l_{j'}\).

\section{Calcolo deduttivo NUOVO CAPITOLO}

Abbiamo definito una logica come una tripla \((\mathcal{L}, \models, \vdash)\), dove \(\mathcal{L}\) è un linguaggio, \(\models\) è la relazione di soddisfacibilità che ne esprime la semantica, e \(\vdash\) è un calcolo, che ci permette di dedurre formule a partire da altre, o verificare se una formula è valida o meno.

Abbiamo già visto una formula di calcolo tramite le tabelle di verità, ma questo tipo di calcolo non \qi{scala}, poiché ad esempio nel prim'ordine non sarà possibile ridursi a un numero finito di casi, come abbiamo visto per il prim'ordine.

Il calcolo che vedremo sarà di pura manipolazione simbolica, noi ne vedremo due, il primo è il \keyword{calcolo di deduzione naturale}.

Un calcolo in generale è definito da un insieme di formule denominato \keyword{assiomi} \(Axi \subseteq PL\),  e un insieme di regole di trasformazione dette \keyword{regole di inferenza} \(R\) che a partire da un numero finito di formule, dette \keyword{premesse}, ci permettono di dedurre una nuova formula, detta \keyword{conclusione}.

Ovviamente non tutte le regole di trasformazioni sono regole d'inferenza valide, poiché vorremmo che queste regole preservassero delle proprietà semantiche. Tramite esse potremmo definire la relazione di \keyword{deducibilità} \(\Gamma \vdash \phi\), se partendo da \(\Gamma\) applicando regole in \(R\) si potrà dedurre \(\phi\). L'obiettivo sarebbe fare in modo in maniera puramente sintattico per ottenere la conseguenza logica, ovvero vorremmo che:
\[\Gamma \models \phi \iff \Gamma \vdash \phi\]

Per quanto riguarda gli assiomi, sono un insieme minimale di tautologie dalle quali tramite il teorema di sostituzione varranno in diverse istansiazioni.\footnote{In un certo senso la regola di sostituzione fa parte di \(R\).}

Il primo calcolo che vedremo però non ha assiomi, ma solo regole di inferenza, e si chiama \keyword{calcolo di deduzione naturale}.

Ovviamente per definirlo è sufficiente definire le regole, il cui numero dipende dal numero di simboli logici che abbiamo. In particolare avremo regole associate ad ogni operatore logico, e verrà usato \(\bot \equiv P\land \neg P\) come simbolo di contraddizione. Le regole catturano il significato semantico degli operatori logici.

\todo{Inserisci regole di deduzione naturale}

AND
\(A \land B \vdash A\) \(A \land B \vdash B\) \(A , B \vdash A\land B\)
IMPLIES
\(A \to B, A \vdash B\) \(B \vdash A \to B\), questa regola mi permette di eliminare un assunzione, poiché se da \(A\) deduco \(B\), allora posso dedurre \(A \to B\) senza più assumere \(A\).
OR
\(A \vdash A \lor B\) \(B \vdash A \lor B\) \([A]\ldots C ,[B]\ldots C, A \lor B \vdash C\) 
NOT
\(\neg A , A \vdash \bot\) \([A]\ldots \bot \vdash \neg A\)
BOT
\(\bot \vdash A\) \([A] \ldots \bot \vdash A\) (\RAA)

Per il processo di deduzione l'idea è che partiremo da delle assunzioni e usando passi atomici di deduzione, che sono applicazioni di regole di inferenza, arriveremo a dedurre la formula che ci interessa. Alcune delle assunzioni potranno essere eliminate, tutte le altre saranno invece assunzioni residue, che andranno a formare \(\Gamma\) nella deducibilità \(\Gamma \vdash \phi\).

\chapter{Lezione 24/03}

Da queste regole dobbiamo costruire delle deduzioni.

Il concetto di deduzione formalmente cambia dipendentemente dal calcolo.

Per questo tipo di calcolo, una deduzione avrà una forma di un albero

\begin{definition}{Deduzione nel calcolo naturale}
  Dato un insieme di formule \(\Gamma\), una \keyword{deduzione} di \(\phi\) da \(\Gamma\) (denotata come \(\Gamma \vdash \phi\) e detta \keyword{sequente}) è un albero che ha come radice \(\phi\), come foglie formule di \(\Gamma\) o assunzioni scaricate, e in cui ogni nodo interno è ottenuto applicando una qualche regola di inferenza del calcolo naturale con i suoi figli come premesse.
\end{definition}

\begin{example}{}
  Supponiamo di voler dimostrare che \(\vdash A \to \neg \neg A\).\footnote{Insichiamo con \(\vdash A\) per indicare che \(A\) è dimostrabile senza ipotesi, ovvero che \(A\) è un teorema.}

  Chiaramente \(A \equiv \neg \neg A\), quindi ci aspettiamo di poter dimostrare questa formula.

  Iniziamo assumendo \(A\) e \(\neg A\), quindi:
  \begin{prooftree}\label{prtree:natural_deduction_example}
    \AxiomC{\({[A]}_2\)}
    \AxiomC{\({[\neg A]}_1\)}
    \BinaryInfC{\(\bot\)}
    \RuleLabel{\(\neg I_1\)}
    \UnaryInfC{\(\neg\neg A\)}
    \RuleLabel{\(\to I_2\)}
    \UnaryInfC{\(A \to \neg\neg A\)}
  \end{prooftree}

  Andiamo a etichettare le inferenze scaricate con i numeri delle regole di inferenza che le scaricano, in modo da poterle identificare facilmente.

  Possiamo vedere la rappresenzazione come albero della deduzione come segue:
  \begin{center}
    \begin{tikzpicture}
      \node (root) {\(A \to \neg\neg A\)};
      \node (n1) [above =of root] {\(\neg\neg A\)};
      \node (n2) [above =of n1] {\(\bot\)};
      \node (n3) [above left=of n2] {\(A\)};
      \node (n4) [above right=of n2] {\(\neg A\)};
      \draw (n1) -- (root);
      \draw (n2) -- (n1);
      \draw (n3) -- (n2);
      \draw (n4) -- (n2);
    \end{tikzpicture}
  \end{center}
\end{example}

Altri esempi. Possiamo dimostrare proprio \(\vdash A \leftrightarrow \neg \neg A\).

Metà lo abbiamo già dimostrato in \Cref{prtree:natural_deduction_example}, ora dobbiamo dimostrare \(\vdash \neg \neg A \to A\).
Infatti avendo una dimostrazione di \(\neg \neg A \to A\) potremmo usare la regola \(\land I\) per ottenere \(\vdash (A \to \neg \neg A) \land (\neg \neg A \to A)\).

Per fare questo possiamo andare a retroso. Per dimostrare un implicazione \(\neg \neg A \to A\) realisticamente useremo la regola \(\to I\), e questo ci regala un assunzione di \(\neg \neg A\) che dobbiamo usare per dimostrare \(A\).

Non abbiamo niente che ci permetta di dimostrare \(A\) direttamente, ma usando \(\RAA\) possiamo farlo aggiungendo l'assunzione scaricata di \(\neg A\) che ci permette di dimostrare \(\bot\). Dobbiamo mostrare quindi che possiamo dimostrare \(\bot\) da \(\neg A\) e \(\neg \neg A\). Ma questo è facile, basta usare la regola \(\neg E\) con \(\neg A\) e \(\neg \neg A\) come premesse.

Proviamo ora che \((P\to Q)\to((Q\to R) \to (P\to R))\).

Applichiamo lo stesso approccio, vogliamo dimostrare quest'implicazione partendo da un insieme vuoto di premesse, quindi dovremo scaricare ogni assunzione che facciamo.

Iniziamo con \(\to I_1\), allora avremo come antecedente \((Q\to R)\to (P\to R)\) aggiungendo l'assunzione scaricable di \(P\to Q\). Ma il nostro nuovo antecedente è un'implicazione, quindi applichiamo \(\to I_2\) e otteniamo come nuovo antecedente \(P\to R\) aggiungendo l'assunzione scaricabile di \(Q\to R\). Ora abbiamo \(P\to R\) da dimostrare, e possimao applicare ancora \(\to I_3\) per ottenere \(R\) come nuovo antecedente aggiungendo l'assunzione scaricabile di \(P\). 

Avendo già come assunzione \(Q \to R\), assumendo \(Q\) potrei dedurre \(R\) tramite \(\to E\). Devo trovare un modo per dedurre \(Q\) o scaricarla. Ma avendo già \(P\to Q\) e \(P\) come assunzioni scaricate posso dedurre \(Q\) tramite \(\to E\).

A differenza del caso precedente, qui le assunzioni scaricate ci guidano nella scelta delle regole.

Proviamo ora a dimostrare \(\neg(P\land\neg P)\).

Il connettivo principale qui è \(\neg\), quindi la strada più promettente è quella di usare \(\neg I_1\), che ci fornisce anche un'assunzione scaricabile di \({[P\land \neg P]}_1\). Dobbiamo allora dimostrare \(\bot\) da \(P\land \neg P\). Per ottenere \(\bot\) è ragionevole pensare di usare \(\neg E\) poiché dalle assunzioni che posso scaricare posso sia dimostrare \(P\) che \(\neg P\) da \(P\land \neg P\) tramite \(\land E_l\) e \(\land E_r\) rispettivamente, e poi usando \(\neg E\) con \(P\) e \(\neg P\) come premesse ottengo \(\bot\).

Qui entrambe le assunioni sono scaricate, poiché per far si che l'albero di deduzione sia ben formato le assunzioni che scarico devono essere in foglie che hanno nel percorso dalla radice a tale foglia almeno una regola di inferenza che le scarica. In questo caso, essendo che \({[P\land\neg P]}_1\) è scaricata dalla regola in radice, hanno entrambe un percorso che va dalla radice a loro che contiene una regola di inferenza che le scarica, quindi sono entrambe scaricate.

In altre parole, per poter essere sicuri che una assunzione è scaricata, dobbiamo assicurarci che tutte le foglie in cui è usata appartengono a sottoalberi di deduzione radicate in regole di inferenza che scaricano tale assunzione.

Proviamo allora a dimostrare \(P \lor \neg P\).

Un idea potrebbe essere quella di usare \(\lor I\) per ottenere \(P \lor \neg P\) da \(P\). Ora però dobbiamo dimostrare \(P\).
L'unica strada ragionevole sembra essere quella di usare \({\RAA}_1\) per ottenere \(P\) da \(\bot\) regalandoci l'assunzione scaricabile di \({[\neg P]}_1\). Ora però dobbiamo dimostrare \(\bot\) da \(\neg P\). Dovremmo però in qualche modo trovare un modo di scaricare \(P\).

Possiamo notare che \(\lor E\) può essere applicata anche a partire da \(\neg P\). Allora possiamo dedurre \(P \lor \neg P\) dall'assunzione scaricata di \(\neg P\). Per arrivare a \(\bot\) ci basta assumere \(\neg (P\lor \neg P)\). Ma se possiamo fare questo posso applicare tale assunzione anche all'attuale radice, ottenendo \(\bot\) da \(P \lor \neg P\) e \(\neg (P\lor \neg P)\) tramite \(\neg E\). Ora però abbiamo ottenuto \(\bot\) con un unica assunzione non scaricata che è \(\neg (P\lor \neg P)\), ma allora per \(\RAA\) posso ottenere \(P \lor \neg P\) scaricando l'assunzione di \(\neg (P\lor \neg P)\).

Questo ci mostra che nonostante la regola di deduzione finale sia \(\RAA\) anche partendo da un altra regola di deduzione (\(\lor I\)) siamo riusciti a dimostrare lo stesso \(P \lor \neg P\).

Abbiamo già accennato che \(P\to Q \equiv \neg P \lor Q\). Di conseguenza ci aspettiamo che si possano dimostrare entrambi i sequinti \(P\to Q \vdash \neg P \lor Q\) e \(\neg P \lor Q \vdash P\to Q\).

In questo caso abbiamo \(\Gamma \neq \emptyset\), in particolare non dovremmo scaricare \(P\to Q\).

Dalla dimostrazione di prima abbiamo visto che un altro modo per ottenere un \(\lor\) è tramite \(\RAA\).
Andiamo quindi ad aggiungere come assunzione scaricata \(\neg(\neq P \lor Q)\), e dobbiamo dimostrare \(\bot\).
Ragionevolmente useremo una \(\neg E\) e ci servono due formule. Da \(P \to Q\) ci aspettiamo che potremo ottenere \(Q\), mentre da \(\neg(\neg P \lor Q)\) ci aspettiamo di poter ottenere \(\neg Q\). Quindi possimao dedurre \(\bot\) proprio da \(Q\) e \(\neg Q\) tramite \(\neg E\). Dobbiamo allora dimostrare effettivamente \(Q\) da \(P \to Q\) e \(\neg Q\) da \(\neg(\neg P \lor Q)\).

Per ottenere \(Q\) possiamo usare \(\to E\) con \(P\) come premessa, ma \(P\) non è un assunzione scaricata, quindi dobbiamo trovare un modo per dimostrarla. Intuitivamente anche questa potrebbe essere ottenuta da \(\neg (\neg P \lor Q)\).

Per l'altro lato, per ottenere \(\neg Q\) possiamo usare \(\neg I\). Dobbiamo allora dimostrare \(\bot\) da \(Q\) scaricata. Ma da \(Q\) per \(\lor I\) possiamo ottenere \(\neg P \lor Q\) che con \(\neg (\neg P \lor Q)\) ci permette di ottenere \(\bot\) tramite \(\neg E\).

Ci manca quindi solo dimostrare \(P\). Possiamo quindi usare \(\RAA\) per ottenere \(P\) da \(\bot\) scaricando l'assunzione di \(\neg P\). Qui ancora per \(\lor I\) possiamo ottenere \(\neg P \lor Q\) da \(\neg P\), e con \(\neg (\neg P \lor Q)\) otteniamo \(\bot\) tramite \(\neg E\).

Proviamo allora a dimostrare ora l'altro verso, ovvero \(\neg P \lor Q \vdash P\to Q\).

Dobbiamo concludere \(P\to Q\), e il tentativo più promettente è quello di usare \(\to I\) per ottenere \(Q\) da \(P\) scaricata. Ora però dobbiamo dimostrare \(Q\) da \(P\) e \(\neg P \lor Q\). Avendo un \(\lor\) posso pensare di usare \(\lor E\) per ottenere \(Q\) da \(\neg P \lor Q\) che mi fornisce due assunzioni scaricabili \(\neg P\) e \(Q\) a patto che da queste io possa dimostrare \(Q\). Da \(Q\) è chiaro che possa dimostrare \(Q\). Per quanto riguarda \(\neg P\) e avendo \(P\) scaricata da \(\to I\) otteniamo \(\bot\), che tramite \(\bot E\) ci permette di dimostrare \(Q\).

\chapter{Lezione 26/03}

Iniziamgo con qualche esercizio.
\section{TEMP}
\subsection*{Dimostrazione di \(P\land Q \vdash \neg (\neg P \lor \neg Q)\).}

Essendo che l'operatore principaele è la una contradizzione possiamo inziare con una \(\neg I\).

\begin{prooftree}
  \AxiomC{\(\bot\)}
  \RuleLabel{\(\neg I_1\)}[\(\neg P \lor \neg Q\)]
  \UnaryInfC{\(\neg (\neg P \lor \neg Q)\)}
\end{prooftree}

Essendo che semanticamente sappiamo che da \(P \land Q\) segue che sia \(P\) che \(Q\) sono veri, quindi è ragionevole pensare di poter ottenere una contraddizione da \(\neg P \lor \neg Q\). Qui ci tocca usare \(\lor E\) poiché per usare una \(\neg E\) dovrei usare ciò che sto cercando di dimostrare, e questo non è possibile.

\begin{prooftree}
  \AxiomC{\(\bot\)}
  \AxiomC{\(\bot\)}
  \AxiomC{\({[\neg P \lor \neg Q]}_1\)}
  \RuleLabel{\(\lor E\)}[\(\neg P\), \(\neg Q\)]
  \TrinaryInfC{\(\bot\)}
  \RuleLabel{\(\neg I_1\)}
  \UnaryInfC{\(\neg (\neg P \lor \neg Q)\)}
\end{prooftree}

Adesso mi serve trovare una contraddizione sia da \(\neg P\) che da \(\neg Q\). Per entrambi sarà sufficiente usare una \(\land E\) per ottenere \(P\) e \(Q\) da \(P \land Q\).

%TODO AGGIORNA
\begin{prooftree}
  \AxiomC{\(\bot\)}
  \AxiomC{\(\bot\)}
  \AxiomC{\({[\neg P \lor \neg Q]}_1\)}
  \RuleLabel{\(\lor E\)}[\(\neg P\), \(\neg Q\)]
  \TrinaryInfC{\(\bot\)}
  \RuleLabel{\(\neg I_1\)}
  \UnaryInfC{\(\neg (\neg P \lor \neg Q)\)}
\end{prooftree}

Per quanto riguarda l'altro lato il ragionamento è diverso. Iniziamo con \(\land E\).

Quindi ci serve dimostrare separatamente \(P\) e \(Q\). Ma qui è facile, basta usare \(\bot\) da \(\RAA\) in entrambi i casi in modo da ottenere come assunzioni scaricabili i negati, da cui posso ottenere per \(\lor I\) la negata di \(\neg (\neg P \lor \neg Q)\) in modo da ottenre la contraddizione che cercavo via \(\neg E\).

Da questo usando \(\to I\) entrambi i lati e \(\land I\) posso ottenre che le due cose sono equivalenti.

\subsection*{Dimostrazione di \((P\to Q) \leftrightarrow (\neg Q \to \neg P)\).}

Andiamo anche qui a dimostrare i due versi. Iniziamo con \((P\to Q) \to (\neg Q \to \neg P)\).
Partiamo con \(\to I\) che im regala l'assunzione \(P\to Q\). Devo dimostare \(\neg Q \to \neg P\). Questa è un altra implicazione, quindi posso usare \(\to I\) di nuovo, e ottenere come assunzione \(\neg Q\). A questo punto devo dimostare \(\neg P\). A questo punto posso usare \(\neg I\) per ottenere una contraddizione da \(P\). Con \(P\) e \(P\to Q\) posso ottenere \(Q\) con \(\to E\), e con \(Q\) e \(\neg Q\) ottengo la contraddizione che cercavo.

Per l'altro verso, partiamo con \(\to I\) che ci regala l'assunzione \(\neg Q \to \neg P\). Devo dimostrare \(P\to Q\). Questa è un altra implicazione, quindi posso usare \(\to I\) di nuovo, e ottenere come assunzione \(P\). A questo punto devo dimostrare \(Q\). A questo punto posso usare \(\neg I\) per ottenere una contraddizione da \(\neg Q\). Con \(\neg Q\) e \(\neg Q \to \neg P\) posso ottenere \(\neg P\) con \(\to E\), e con \(P\) e \(\neg P\) ottengo la contraddizione che cercavo.

\subsection{Dimostrazione di \((\neg \neg (A \to B)) \leftrightarrow (A \to B)\)}

Questa ha la stessa struttura di \(\neg \neg A \leftrightarrow A\), che abbiamo già dimostrato, ma allora possiamo prendere l'albero di derivazione di \(\neg \neg A \leftrightarrow A\) e sostituire \(A\) con \(A \to B\) in modo da ottenere la dimostrazione che cerchiamo.

\subsection*{Dimostrazione di \((A \to B) \leftrightarrow (\neg \neg A \to \neg \neg B)\)}

Anche qua conviene andare per passi. Iniziamo con \((A \to B) \to (\neg \neg A \to \neg \neg B)\). Usa \(\to I\), \(\neg I\) e poi da \(\neg \neg A\) sappiamo già che equivale ad \(A\), quindi possiamo usare \(\to E\) per ottenere \(B\), e da \(B\) e \(\neg \neg B\) otteniamo la contraddizione che cercavamo.

Altro lato come esercizio, ma è simile al precedente.

\subsection*{Dimostrazione di \((A \land B)\lor C \vdash (A \lor C) \land (B \lor C)\).}

Iniziamo con \(\land I\). Dobbiamo quindi dedurre i due congiunti. Per usare l'assunzione dovremo prima o poi usare una \(\lor E\). Le due assunzioni che \(\lor E\) scarica sono i disgiunti quindi potrò usare anche quelle.

Banale falla a casa.

Facciamo l'altro versio. Iniziamo con \(\lor I\), devo dimostrare \((A \land B)\) oppure \(C\). Per dimostrare \(A \land B\) devo usare \(\land I\), quindi devo dimostrare \(A\) e \(B\). Dall'assunzione che non devo scaricare posso ottenere \(A \lor C\) e \(B \lor C\) con \(\land E\) che potrò usare in un \(\lor E\). Allora posso usare \(\lor E\) con \(A \lor C\) per ottenere \((A \land B) \lor C\) scaricando la \(A\). Per scaricare la \(B\) lo faccio ancora.

Questo è un altro esempio in cui dobbiamo aggiungere cose sotto l'ipotizzata radice, non è detto che riusciamo sempre a imbroccare la regola giusta alla prima botta.

\subsection*{Dimostrazione di \((A \lor B) \land C \vdash (A \land C) \lor (B \land C)\).}

\(\lor I\), devo dimostrare \(A \land C\).
\(\land I\) devo dimostrare \(A\) e \(C\).
\(C\) è da \(\land E\) di \((A \lor B)\land C\). Per \(A\) devo usare \(\lor E\) con \((A \lor B)\). Per la prima dimostrazione di \(A\) è gratis da assunzione di \(A\).
Ora devo dimostrare \(A\) assumendo \(B\).

Continua a casa.

Esercizi

\[\neg (A \to B), (\neg B) \to C \vdash C\]
\[((A\land B)\lor (\neg C \lor B))\to (C \to B)\]
\[(A \to (B\lor C))\lor((\neg A \land \neg C)\to B)\]

\section{Relazione fra deduzione e conseguenza logica}

\begin{theorem}{Monotonicità di \(\vdash\)}
  \(\Gamma \vdash \phi \implies \Gamma \cup \{\psi\} \vdash \phi\)
\end{theorem}

Intuitivamente questo è vero anche se introduco contraddizioni, grazie alla regola \(\bot E\).

\begin{note}{}
  Questa proprietà di generalizza in \(\Gamma \vdash \phi \implies \Gamma \cup \Gamma' \vdash \phi\).
\end{note}

\begin{proof}
  Per definizione di deduzione, \(\Gamma \vdash \phi\) se e solo se esiste un albero di derivazione ben formato radice etichettata con \(\phi\) e le cui foglie o sono etichettate con formule di \(\Gamma\) o sono assunzioni scaricabili. 

  In tale definizione non c'è alcun obbligo di usare ogni formula di \(\Gamma\) come etichetta di una foglia. Pertanto, aggiungendo all'insieme delle premesse \(\Gamma\) una nuova formula \(\psi\) non usata nell'albero di derivazione, lo stesso identico albero di derivazione rimane valido per dimostrare \(\Gamma \cup \{\psi\} \vdash \phi\).
\end{proof}

\begin{note}{}
  NON È VERO CAMBIA.
  Questa proprietà vale nella logica classica, ma non in tutte le logiche. Ad esempio, nella logica intuizionista, in cui vogliamo ottenere deduzioni costruttive, l'aggiunta di una nuova formula \(\psi\) potrebbe non essere innocua, poiché potrebbe portare a contraddizioni.
\end{note}

Possiamo vedere un analogo del \Namedcref{thm:semantic_deduction}. Infatti questa proprietà è vera anche per \(\vdash\).

\begin{theorem}{Teorema di Deduzione}
  \(\Gamma, \psi \vdash \phi \iff \Gamma \vdash \psi \to \phi\)
\end{theorem}

\begin{proof}
  Andiamo a dimostrare entrambi i versi. Iniziamo con \(\Gamma, \psi \vdash \phi \implies \Gamma \vdash \psi \to \phi\). Per definizione di deduzione, esiste un albero che dimostra \(\Gamma, \psi \vdash \phi\). Ma allora mi basta applicare la regola \(\to I\) a questo albero, scaricando l'assunzione \(\psi\) e ottenendo così un albero che dimostra \(\Gamma \vdash \psi \to \phi\).

  L'altro verso, \(\Gamma \vdash \psi \to \phi \implies \Gamma, \psi \vdash \phi\), è ancora più semplice. Per definizione di deduzione, esiste un albero che dimostra \(\Gamma \vdash \psi \to \phi\). A questo punto, aggiungendo l'assunzione \(\psi\) all'albero, possiamo applicare la regola \(\to E\) per ottenere un albero che dimostra \(\Gamma, \psi \vdash \phi\).
\end{proof}

Queste cose mi iniziano a suggerire che le due relazioni godono delle stesse proprietà, e che quindi forse sono equivalenti.

\begin{theorem}{Regola del taglio}
  Se \(\Gamma, \psi \vdash \phi\) e \(\Gamma' \vdash \psi\) allora \(\Gamma \cup \Gamma' \vdash \phi\)
\end{theorem}

\begin{proof}
  Intuitivamente basta sostituire a ogni foglia etichettata con \(\psi\) nell'albero di derivazione che dimostra \(\Gamma, \psi \vdash \phi\) l'albero di derivazione che dimostra \(\Gamma' \vdash \psi\). In questo modo otteniamo un albero di derivazione che dimostra \(\Gamma \cup \Gamma' \vdash \phi\).

  Un altro modo per dimostrare questo è utilizzare la monotonicità e il teorema di deduzione.
  Infatti per monotonicità abbiamo sia \(\Gamma, \Gamma', \psi \vdash \phi\) e \(\Gamma \cup \Gamma'\vdash \psi\).
  A questo punto, usando il teorema di deduzione, otteniamo \(\Gamma, \Gamma' \vdash \psi \to \phi\). 

  Posso quindi costruire un albero con \(\phi\) in radice che usa come regola in radice \(\to E\) e come sottoalberi l'albero che deriva \(\psi\) e l'albero che deriva \(\psi \to \phi\). In questo modo otteniamo un albero di derivazione che dimostra \(\Gamma \cup \Gamma' \vdash \phi\).
\end{proof}

\begin{definition}{Inconsistenza}
  Un insieme di formule \(\Gamma\) è inconsistente se \(\Gamma \vdash \bot\).

  Una definizione equivalente è che \(\Gamma\) è inconsistente se \(\Gamma \vdash \phi\), \(\forall \phi \in \PL\).
\end{definition}

L'inconsistenza è la versione sintattica all'interno del calcolo dell'essere insoddisfacibile.

La prossima volta dimostreremo correttezza e completezza del calcolo.

Ovviamente non si bastano singolarmente, perché per garantire la sola correttezza basta usare come relazione di deduzione una relazione vuota, mentre per garantire la sola completezza basta usare come relazione di deduzione una relazione che contiene tutte le coppie \((\Gamma, \phi)\). Quindi averle da sole non è particolaremente interessante.
